\chapter*{Introducción}
\addcontentsline{toc}{chapter}{Introducción}

Este libro es el resultado de la experiencia de los autores al orientar la asignatura de Cálculo Integral en la Universidad Distrital Francisco José de Caldas durante varios años. Bajo esta premisa, el publico objetivo del presente texto son los estudiantes de tecnología e ingeniería que tomen la asignatura de Cálculo Integral, con la intención de que puedan usarlo como guía durante el desarrollo del curso.
En este libro, se desarrolla la teoría de un curso típico de cálculo integral de una variable real; por lo que se realizan las demostraciones de algunos teoremas, y, de los que no, se hace una explicación gráfica. Además, se presenta un abanico de variados ejemplos referentes a distintos temas con una explicación cuidadosa; en algunos se incluye un mayor grado de dificultad en comparación con los ejemplos que normalmente se encuentran en los textos clásicos de cálculo integral. Seguramente con esta dinámica le permitirá al estudiante desarrollar los ejercicios de este y de cualquier otro libro.

Ahora, el presente documento está conformado por cinco capítulos. En el capítulo uno se desarrolla el concepto de integral definida. En el capítulo dos se presentan las integrales indefinidas y el teorema fundamental del cálculo en sus partes I y II, con sus implicaciones. En el capítulo tres se presentan diversos métodos de integración. En el capítulo cuatro se tratan las aplicaciones de la integral definida. Y en el capítulo cinco se desarrollan las integrales impropias y la integración numérica.

El libro se escribió buscando incluir novedades en comparación con los libros de texto que se encuentran en la literatura. Por ejemplo, se incluye el teorema (1.2) y se demuestra usando inducción matemática, con el fin de probar la integral de $x$ a la $m$.
Por otro lado, también se incluyen integrales definidas de funciones inversas mediante sumas o restas de áreas y sin usar el método de integración por partes. Se presenta el teorema (1.8), que indica cómo calcular esa integral para funciones continuas y crecientes.

Como aporte didáctico, este libro cuenta con algunos enlaces a \textit{GeoGebra}, donde el lector puede interactuar, por ejemplo, con la construcción de los sólidos de revolución que se involucran en el cálculo de un volumen usando los métodos de discos, arandelas, capas o casquillos cilíndricos. También se incluyen animaciones de volumen de sólidos con función de área transversal conocida. Esta herramienta estimulará el trabajo del estudiante y le ayudará a visualizar su proceso y a comprender mejor las magnitudes clave para el planteamiento de las integrales que llevan a la medición de un volumen o de un área superficial. 

En efecto, esto mejorará el aprendizaje de los temas planteados y apoyará el trabajo en el aula desarrollado por el docente que imparte el curso. Todas las gráficas que se muestran en este libro fueron construidas en \textit{GeoGebra} por los autores.
